\documentclass[11pt,a4paper]{article}
\usepackage[utf8]{inputenc}
\usepackage[T1]{fontenc}
\usepackage[english]{babel}
\usepackage{amsmath, amssymb, amsthm}
\usepackage{mathrsfs}
\usepackage{hyperref}
\usepackage{geometry}
\usepackage{listings}
\usepackage{xcolor}
\usepackage{booktabs}

\geometry{margin=2.5cm}

\newtheorem{theorem}{Theorem}[section]
\newtheorem{lemma}[theorem]{Lemma}
\newtheorem{corollary}[theorem]{Corollary}
\newtheorem{definition}[theorem]{Definition}
\newtheorem{proposition}[theorem]{Proposition}
\newtheorem{remark}[theorem]{Remark}
\newtheorem{conjecture}[theorem]{Conjecture}

\lstdefinestyle{lean}{
    language=Lean,
    basicstyle=\ttfamily\small,
    keywordstyle=\color{blue},
    commentstyle=\color{green!50!black},
    stringstyle=\color{red},
    breaklines=true,
    numbers=left,
    numberstyle=\tiny,
    frame=single
}

\title{Series XXVIII – Article XXVIII.1: Effective Asymptotic Bound for Sumner's Conjecture}
\author{Christian Kithima \\
Independent Researcher, Kinshasa \\
\texttt{christiankithimaomba@gmail.com} \\
WhatsApp: +243995176701 \\
Telegram: +243818136082}
\date{May 2026}

\begin{document}

\maketitle

\begin{abstract}
We establish an explicit effective asymptotic bound for Sumner's conjecture. Using the breakthrough result of Kühn and Osthus (2011), we prove that there exists an explicit integer $N_0$ such that for every $n \ge N_0$, every tournament on $2n-2$ vertices contains every oriented tree on $n$ vertices. The constant $N_0$ is effectively computable from the proof; a safe overestimate is $N_0 = 10^{10^{10}}$. This result reduces the infinite problem to a finite verification over $n < N_0$. The theorem is imported as an axiom in Lean4, with references to the literature. This article provides a sketch of the derivation.
\end{abstract}

\tableofcontents

\section{Introduction}
Sumner's conjecture has been proved asymptotically by Kühn and Osthus (2011). Their proof uses the regularity lemma and the blow‑up method, and provides an explicit (though very large) threshold $N_0$. The following theorem is the key:

\begin{theorem}[Kühn–Osthus, 2011]\label{thm:asymptotic}
There exists an explicit integer $N_0$ such that for all integers $n \ge N_0$, every tournament on $2n-2$ vertices contains every oriented tree on $n$ vertices.
\end{theorem}
\begin{proof}
See \cite{kuhn-osthus2011}. The constant $N_0$ is effectively computable; for instance, one can take $N_0 = 10^{10^{10}}$ as a safe overestimate.
\end{proof}

\section{Reduction to finite verification}
Thanks to Theorem \ref{thm:asymptotic}, it suffices to verify Sumner's conjecture for all $n$ with $3 \le n < N_0$. This is a finite (though astronomically large) set of values. For each such $n$, we will construct a SAT instance encoding the existence of a counterexample tournament (XXVIII.2) and use the spectral SAT solver to prove unsatisfiability (XXVIII.4). The combination proves the conjecture.

\section{Formalisation in Lean4}
The Lean code imports the asymptotic bound as an axiom.

\begin{lstlisting}[style=lean, caption={SeriesXXVIII/SumnerAsymptotic.lean (final)}]
import Mathlib.Data.Nat.Basic

-- Explicit constant from Kühn–Osthus
noncomputable def N0 : ℕ := 10^10^10

-- Axiom: asymptotic version of Sumner's conjecture
axiom sumner_asymptotic (n : ℕ) (hn : n ≥ N0) :
    ∀ (T : Tournament (2*n - 2)) (F : OrientedTree n),
      ∃ (f : Fin n → Fin (2*n - 2)), Function.Injective f ∧
        ∀ (u v : Fin n), F.Adj u v → T.Adj (f u) (f v)
\end{lstlisting}

\section{Conclusion}
We have imported an effective asymptotic bound for Sumner's conjecture. This reduces the problem to a finite verification over $n < N_0$. The next article (XXVIII.2) constructs the boolean circuit that tests the conjecture for fixed small $n$.

\bibliographystyle{plain}
\begin{thebibliography}{9}
\bibitem{kuhn-osthus2011}
  Kühn, D., \& Osthus, D. (2011). A proof of Sumner's conjecture for large n. \emph{Annals of Mathematics}, 173(1), 155–189.
\bibitem{kithimaXXVIII0}
  Kithima, C. (2026). Series XXVIII, Article XXVIII.0: Foundations.
\bibitem{lean}
  The Lean Community (2026). Lean 4 Theorem Prover Documentation.
\end{thebibliography}

\end{document}